% ULM final-project report outline. Replace drafting notes with scientific prose.
% Official ACL 2023 files, included unchanged alongside this document:
% https://2023.aclweb.org/downloads/acl2023.zip
% Build from docs/: latexmk -pdf -interaction=nonstopmode report-outline.tex
% Alternatively, use pdflatex directly for the outline. After enabling the
% bibliography, run pdflatex, bibtex, pdflatex, pdflatex in that order.
% Course requirements override generic ACL conference submission rules:
% abstract <=150 words; required section order; AI declaration after References.
% No course page limit or institutional declaration wording was supplied.
% Optional Conclusion and Acknowledgments are included and may be removed.
% Writing order suggested by the course: notes for Introduction, then Methods,
% Results, Introduction, Discussion/Conclusion, and finally Abstract.
\pdfoutput=1
\documentclass[11pt]{article}
\usepackage[final]{acl2023}
\usepackage{times}
\usepackage{latexsym}
\usepackage[T1]{fontenc}
\usepackage[utf8]{inputenc}
\usepackage{microtype}
\usepackage{inconsolata}
\usepackage{amsmath,amssymb}
\usepackage{booktabs,tabularx}
\usepackage{xurl}

\newcommand{\todo}[1]{\textit{[TODO: #1]}}
\newcommand{\sourcefile}[1]{\par\noindent\textit{Repository evidence:} \path{#1}.\par}
\newcommand{\placeholder}{\textit{TBD}}
\newif\ifreferencesready
% Uncomment after adding and citing real entries in report-references.bib:
% \referencesreadytrue

\title{PassGame: Preserving Accumulated Constraints in Multi-Turn Language Model Evaluation}
\author{\textbf{Author name(s)} \\
  Department / University \\
  ULM course, semester / year \\
  \texttt{email@example.com}}
% Add student ID(s), instructor, and submission date only if required by the
% course. Keep the author block within the template's allocated title area.

\begin{document}
\maketitle

\begin{abstract}
\todo{Write this last, in no more than 150 words. State the problem and research
question; summarize the benchmark, models, and experimental conditions; give
the principal quantitative finding; and explicitly state the study's conclusion.
Replace this guidance with a single self-contained paragraph.}
\end{abstract}

\section{Introduction}
\label{sec:introduction}

\subsection{Motivation and Research Context}
\begin{itemize}
  \item Motivate the need to preserve earlier requirements when an assistant
  updates an artifact over several turns. Connect this to instruction following,
  constraint satisfaction, and the relevant ULM class topics.
  \item Introduce PassGame: each turn adds one rule to a persistent artificial
  string, and a deterministic checker evaluates all active rules.
  \item Explain why satisfying a new instruction can conflict with preserving
  old ones. Distinguish this question from following a single isolated instruction.
\end{itemize}

\subsection{Relevant Prior Work and Research Gap}
\begin{itemize}
  \item Discuss the most relevant work from the project proposal and course
  materials on instruction following, multi-turn evaluation, and compositional
  constraints; an exhaustive review is unnecessary.
  \item Explain the specific gap addressed by progressive rule disclosure,
  exact checking, and certificates of local repair interference. Support any
  novelty claim with the cited literature.
  \item Add verified references to \path{report-references.bib}; use
  \verb|\citet{key}| for textual and \verb|\citep{key}| for parenthetical citations.
\end{itemize}

\subsection{Research Questions and Hypotheses}
\begin{itemize}
  \item \textbf{RQ1 / H1:} How does full active-rule success vary with certified
  interference in the model's actual incoming state? State the proposed direction
  and distinguish a descriptive association from a causal effect.
  \item \textbf{RQ2 / H2:} How often does a model satisfy the new rule while
  breaking a previously satisfied rule? Specify unconditional and conditional
  regression measures.
  \item \textbf{RQ3 / H3:} For an identical final rule set and initial string,
  does original versus reverse rule order change whole-game success or the
  number of successfully completed turns? Treat random orders as secondary.
  \item Identify which questions the reported experiments actually address.
  Label any further reminder, parser, or thinking-mode comparison exploratory.
\end{itemize}

\subsection{Contribution and Scope}
\todo{Summarize the implemented benchmark, validation, and empirical study.
Distinguish development-pilot evidence from held-out evaluation and diagnostic
algorithm checks from measurements of language models.}

\section{Methods}
\label{sec:methods}
% Describe the procedure reproducibly. Do not put observed results or their
% interpretation here; those belong in Results and Discussion.

\subsection{Task, String Domain, and Rule Semantics}
\begin{itemize}
  \item Define the incoming string $x_t$, newly revealed rule $r_t$, old rules
  $R_{<t}$, and model response $y_t$. Success requires $y_t$ to satisfy every
  rule in $R_{<t}\cup\{r_t\}$.
  \item Specify strings of length 1--32 over ASCII letters, digits, and
  \verb|!_-()[]|; case sensitivity, one-based positions, repeated digits,
  overlapping substring counts, and comparison operators.
  \item Describe the rule kinds used in the selected dataset. Distinguish
  implemented rule support from actual dataset coverage, including the deferred
  bracket-specific family.
\end{itemize}
\sourcefile{src/passgame/constraints.py}

\subsection{Dataset Construction and Splits}
\begin{itemize}
  \item Describe the v3 neutral-seed progressive suite: eight scheduled rules
  per game; 24 development and 120 held-out original games; six families with
  four development and twenty held-out games each.
  \item Explain the six families: numeric constraints; character counts with
  positions; required substrings with positions; forbidden substrings with
  positions; and numeric interference involving two or three old rules.
  Give a compact worked example.
  \item Specify the generation seed, split policy, exact dataset hashes, and
  final-rule-set separation. Order variants stay in the same split.
  \item Explain private reference paths as satisfiability proofs for every
  cumulative prefix. They are neither model outputs nor fixed interference
  labels for autonomous trajectories.
  \item Describe automated checks with an independent checker and the human
  audit procedure. Report the audit status applicable to the evaluated snapshot;
  the current v3 manifest records it as pending for held-out model evaluation.
\end{itemize}
\sourcefile{data/sequential-v3/manifest.json}
\sourcefile{src/passgame/sequential.py}
\sourcefile{docs/audit-instructions.md}

\subsection{Sequential Interaction and Response Protocol}
\begin{enumerate}
  \item Initialize once from the episode's neutral seed, shared across models.
  \item Present the current string, exactly one new rule, and, in the main
  condition, the exact earlier-rule reminders. Retain the real conversation.
  \item Request a JSON object containing only the \texttt{password} field,
  whose value is a string. The response schema is shown below.
  \item Parse the response and check the domain and every active rule.
  \item Carry forward a fully successful output; terminate at the first semantic,
  format, or domain failure. Retain the failed attempt.
\end{enumerate}
\noindent\verb|{"password":"..."}|

Describe exact system/user prompts, transport retry limits, and the treatment of
infrastructure failures as missing evidence. State that there are no semantic
retries, oracle repairs, future-rule disclosures, tools, or provider-enforced JSON
mode. Full rewrites are allowed; minimum edit distance is not an output constraint.
Specify strict parsing or the optional single outer Markdown-fence policy for
each condition, without pooling the two protocols.
\sourcefile{src/passgame/protocol.py}
\sourcefile{src/passgame/run.py}

\subsection{Interference Definition and Certification}
\begin{itemize}
  \item Define unit-cost Levenshtein distance and the minimum distances to a
  new-rule-valid string ($d_{\mathrm{new}}$) and an all-rules-valid string
  ($d_{\mathrm{full}}$).
  \item Define interference degree as the minimum number of distinct old rules
  broken across all shortest new-only repairs. Rule IDs may differ across repairs.
  \item Specify the certified labels: no-op ($0,0$), independent ($1,1$),
  interfering ($1,2$), and unknown when certification is inconclusive.
  \item Describe exhaustive domain-valid one-edit neighborhoods generated by
  insertion, deletion, or substitution, the verified two-edit witness, and the bounded
  witness-search settings used by the evaluated source snapshot.
  \item Recompute certificates from actual model states. Describe unknown/no-op
  coverage explicitly and keep legacy controlled pairs separate from v3 episodes.
\end{itemize}
\sourcefile{src/passgame/oracle.py}

\subsection{Models, Conditions, and Execution Settings}
\todo{For every evaluated condition, record the exact model/provider ID and
weight revision; dense/MoE architecture and parameter counts; precision or
quantization; verified thinking mode and reasoning parser; chat template;
temperature, sampling and context limits; total completion-token cap including
reasoning; final-answer extraction; serving engine/version; GPU type/count;
run date; and configuration/output identifiers.}

Describe the chosen track and split, rule horizon, reminder setting, parser
policy, order variants, and any repeated runs or seeds. For v3 order development,
24 final sets times four orders gives 96 games and at most 768 requests per model;
the seeded development suite permits at most 192 requests. Counts are ceilings
because failed episodes stop early. Distinguish configured conditions from
conditions actually evaluated, and document confounds between model settings.
\sourcefile{configs/dev-qwen2.5-7b-v3-order.yaml}
\sourcefile{docs/experiment-guide.md}

\subsection{Outcome Measures and Statistical Analysis}
\begin{itemize}
  \item Define full success, new-rule success, all-old success, retention of
  previously satisfied rules, fractional retention, and newly broken-rule counts.
  Report the eligible denominator for every measure.
  \item Define new-rule success with regression over all retention-eligible
  responses, and regression conditional on new-rule success.
  \item Define whole-game success, successful turns before first failure,
  first-failure stage, reached-stage counts, and survival with at-risk denominators.
  Observed turns include a terminal failed turn; successfully completed turns do not.
  \item Explain invalid-output scoring, unknown broken-rule counts, infrastructure
  missingness, and the exclusion of unattempted later turns from response metrics.
  \item For H3, pair original-minus-reverse outcomes by final rule set, initial
  string, and horizon. Specify 95\% paired-bootstrap intervals, 2,000 seeded
  resamples or the recorded alternative, and missing-block handling. Treat random
  orders as secondary comparisons.
  \item Keep autonomous interference strata descriptive. Do not treat turns as
  independent observations or interpret null/zero-block contrasts as zero effects.
\end{itemize}
\sourcefile{src/passgame/analyze.py}
\sourcefile{docs/preregistration.md}

\subsection{Implementation and Reproducibility}
\todo{Identify Python and dependency versions, source commit, dataset/configuration
hashes, prompts, adapter, frozen run snapshots, and analysis version. Describe
relevant correctness tests and caching/resumption safeguards. Record API tokens,
costs, and GPU runtime where measured. Provide a shareable code/artifact location
and exact reproduction commands; omit credentials and private cluster paths.}
Explain which protocol amendments followed pilot evidence, using their actual
chronology. Do not describe retrospective changes as pre-pilot preregistration.
\sourcefile{docs/experiment-guide.md}
\sourcefile{tests/test_protocol.py}

\section{Results}
\label{sec:results}
% Populate only from a named frozen run and its saved evidence. Do not infer
% completion from a configuration or substitute diagnostics for model results.

\subsection{Run Coverage and Overall Performance}
\todo{Identify the analyzed runs and split. Give planned/completed/incomplete
games, attempted requests, skipped-after-failure turns, and infrastructure
missingness. A complete run can contain failed games. Present full/new/retention
scores, whole-game success, and mean successful turns with their denominators.}

% Keep this placeholder within Results rather than floating into Methods.
\begin{center}
\begin{minipage}{\columnwidth}
\centering
\begin{tabularx}{\columnwidth}{@{}Xrrr@{}}
\toprule
Condition & Games & Success & Turns \\
\midrule
\placeholder & \placeholder & \placeholder & \placeholder \\
\placeholder & \placeholder & \placeholder & \placeholder \\
\bottomrule
\end{tabularx}
\captionof{table}{Results placeholder. Identify each model and protocol condition;
report whole-game success as a count/denominator and mean successful turns
out of eight. Separate seeded/order and development/held-out evidence.}
\label{tab:overall}
\end{minipage}
\end{center}
\sourcefile{outputs/<run>/completion.json}
\sourcefile{outputs/<run>/analysis.json}

\subsection{Interference and Regression: RQ1 and RQ2}
\todo{Show actual-state no-op, independent, interfering-by-degree, and unknown
counts. Report success and regression measures with eligible denominators and
family/stage context. State whether sufficient interfering turns were reached
to address the question. Interpret the observed patterns relative to H1/H2,
without inferring a causal effect from autonomous state strata.}
\sourcefile{outputs/<run>/interference-summary.csv}

\subsection{Rule Order and Survival: RQ3}
\todo{Present per-order whole-game success, mean successful turns, failure-stage
counts, and survival curves with at-risk denominators. Report paired original
minus reverse estimates and 95\% intervals with block counts; show random-order
comparisons as secondary. Interpret the findings relative to H3.}
\sourcefile{outputs/<run>/order-summary.csv}
\sourcefile{outputs/<run>/order-effects.csv}
\sourcefile{outputs/<run>/order-survival.csv}

\subsection{Error Analysis and Additional Conditions}
\todo{Separate new-rule failure, old-rule regression, format/domain failure,
and infrastructure missingness. Include anonymized real response examples with
checker outcomes. Report any reminder, parser, or verified thinking-mode
comparison only when supported by separate compatible runs. Label historical
v2/reset evidence and post-pilot analyses explicitly.}
\sourcefile{outputs/<run>/records.jsonl}

\section{Discussion}
\label{sec:discussion}

\subsection{Answering the Research Questions}
\todo{Connect the principal findings to the initial questions and related work.
Explain what the measurements support, where hypotheses were unsupported or
unresolved, and why the results matter for persistent constraint following.
Distinguish observed errors from speculative explanations about model reasoning.}

\subsection{What Worked and What Did Not}
\todo{Reflect on benchmark construction, exact checking, progressive starts,
interference exposure, rule-order comparisons, formatting, and token limits.
Explain how pilot observations motivated documented changes and how this affects
the strength of the conclusions.}

\subsection{Limitations}
\begin{itemize}
  \item Small, synthetic, template-based data; limited natural-language/artifact
  realism; no claim of generalization to unseen templates or real password security.
  \item Development-driven design changes, pending held-out audit where applicable,
  and limited precision for subgroup or model comparisons.
  \item Model-dependent trajectories, no-op exposure, unknown certificates, and
  survival selection after first-failure termination.
  \item Differences in model size, architecture, precision, thinking mode,
  serving stack, and completion budget; avoid attributing confounded differences
  to one factor.
\end{itemize}

\subsection{Future Work and Broader Considerations}
\todo{Propose follow-up studies motivated by the findings: signed held-out
evaluation, larger or new families, controlled interference exposure, repeated
runs, or richer artifacts. Briefly discuss reproducible sharing, privacy of
execution logs, and measured compute/resource costs where relevant.}

\section{Conclusion}
% Optional: merge into Discussion if a separate conclusion repeats it.
\todo{Give a concise answer to the research question, the strongest supported
finding, and one practical implication or next step. Introduce no new results.}

\section*{Acknowledgments}
% Optional: delete if not applicable.
\todo{Name people, teaching support, and compute/service providers whose help
benefited the project. Verify contributions and avoid implying endorsement.}

% References must precede the AI usage declaration, per the course requirements.
% The empty bibliography is intentional. Enable \referencesreadytrue above
% after adding and citing verified entries; BibTeX then creates this section.
% No fabricated papers or placeholder bibliographic records are supplied.
\bibliographystyle{acl_natbib}
\ifreferencesready
  \bibliography{report-references}
\else
  \section*{References}
  \todo{Add verified bibliography entries, cite them in the paper,
  and enable the bibliography switch in the preamble. The included ACL BibTeX
  style formats the reference list and author--year citations.}
\fi

\section*{AI Usage Declaration}
\label{sec:ai-declaration}
\todo{Disclose all AI assistance used in the project and write-up. The confirmed
use for this outline is Codex assistance to inspect repository documentation,
organize the report, and prepare/check the LaTeX template. Expand this statement
to cover earlier and later project work accurately.}
\begin{itemize}
  \item \textbf{Tools:} tool/provider and model/version if known, with dates or
  project phases; do not guess a model identifier.
  \item \textbf{Uses and scope:} planning, literature search, code generation or
  debugging, experiment setup, analysis, translation/editing, and report writing;
  identify which actually occurred and which artifacts were affected.
  \item \textbf{Human contribution and checking:} who designed, implemented,
  verified, interpreted, and approved the work; how generated material and
  references were checked or revised.
  \item \textbf{Evaluated models:} distinguish models serving as experimental
  subjects from AI tools assisting the authors. Record any further disclosure
  details required by the course's full AI policy.
\end{itemize}

% Optional appendices may be added after the AI declaration if permitted by the
% course: full prompts, rule examples, configurations, reproduction commands,
% additional tables, and an amendment timeline. Keep essential Methods in the
% main paper so that it remains understandable and reproducible on its own.
% Before submission: replace all TODO/TBD text and path placeholders; add real
% citations; verify the abstract word count, course page limit, and AI policy;
% confirm the final section order and compile with no unresolved references.

\end{document}
